\documentclass[11pt]{article}
\usepackage[margin=0.95in]{geometry}
\usepackage{amsmath,amssymb,amsthm}
\usepackage{mathtools}
\usepackage{booktabs}
\usepackage{enumitem}
\usepackage{hyperref}

\hypersetup{colorlinks=true,linkcolor=blue,urlcolor=blue,citecolor=blue}
\setlist[itemize]{leftmargin=1.5em,itemsep=2pt}
\setlist[enumerate]{leftmargin=1.7em,itemsep=2pt}

\newcommand{\R}{\mathbb{R}}
\newcommand{\E}{\mathbb{E}}
\newcommand{\norm}[1]{\left\lVert #1 \right\rVert}
\newcommand{\inner}[2]{\left\langle #1, #2 \right\rangle}
\newcommand{\bigO}{\mathcal{O}}

\newtheorem{proposition}{Proposition}
\newtheorem{corollary}{Corollary}
\newtheorem{definition}{Definition}
\newtheorem{remark}{Remark}

\title{Supporting Theory for Incremental Modality Acquisition\\
\large Frozen-LM Residual Adapters, Interference, and Layer Placement}
\date{March 17, 2026}

\begin{document}
\maketitle

\begin{abstract}
This note collects only the theory needed to support the NeurIPS paper on
incremental modality addition to a frozen language model. The central points
are: (i) gated residual adapters preserve the base model exactly when the added
modalities are inactive, (ii) composition error arises from second-order
cross-terms between modality perturbations, (iii) the empirical synergy metric
is a score-level proxy for this interaction, and (iv) staggered layer placement
removes direct same-layer coupling and replaces it with an attenuated
Jacobian-mediated interaction. We deliberately omit sidecar mixing, unrelated
training objectives, and other speculative extensions.
\end{abstract}

\section{Setup}

Let a frozen decoder-only transformer have layers
\[
h_{\ell+1} = f_\ell(h_\ell), \qquad \ell = 0, \dots, L-1,
\]
with text-conditioned initial state $h_0 = \mathrm{Embed}(x_{\text{text}})$.

For each added modality $m$:
\begin{itemize}
\item a frozen encoder $E_m$,
\item a trainable projector $P_m$,
\item trainable fusion blocks $\Phi_{m,\ell}$ at selected layers $\ell \in S_m$,
\item a gate $g_{m,\ell}$ controlling injection strength.
\end{itemize}

The projected modality representation is
\[
z_m = P_m(E_m(x_m)).
\]

At a fusion layer $\ell$, the model applies pre-activation residual fusion
\begin{equation}
\tilde h_\ell
=
h_\ell
 \sum_{m \in \mathcal{M}_\ell(x)}
 g_{m,\ell}(x)\,\Delta_{m,\ell}(h_\ell, z_m),
\qquad
h_{\ell+1} = f_\ell(\tilde h_\ell),
\label{eq:fusion}
\end{equation}
where $\Delta_{m,\ell} := \Phi_{m,\ell}(h_\ell, z_m)$ and
$\mathcal{M}_\ell(x)$ is the set of active modalities on example $x$.

\paragraph{Scope.}
All claims below concern the frozen-backbone residual-adapter setting in
Eq.~\eqref{eq:fusion}. They do \emph{not} claim that arbitrary fine-tuning
methods preserve the base model, nor that any observed metric always admits an
exact second-order explanation globally.

\section{Exact Gated-Off Invariance}

\begin{definition}[Inactive modality]
A modality $m$ is inactive on input $x$ if its fusion contribution is exactly
zero at every fusion layer:
\[
g_{m,\ell}(x)\,\Delta_{m,\ell}(h_\ell, z_m) = 0
\qquad \forall \ell \in S_m.
\]
This covers the cases ``modality absent'', ``gate forced to zero'', or any
implementation that produces the zero residual when the modality is disabled.
\end{definition}

\begin{proposition}[Exact invariance of the frozen hub]
\label{prop:invariance}
If every added modality is inactive on input $x$, then the model with adapters
produces exactly the same hidden states and outputs as the original frozen
text-only backbone on $x$.
\end{proposition}

\begin{proof}
Under inactivity, the sum in Eq.~\eqref{eq:fusion} is identically zero, so
\[
\tilde h_\ell = h_\ell
\qquad \forall \ell.
\]
Therefore the recurrence reduces to
\[
h_{\ell+1} = f_\ell(h_\ell),
\]
which is exactly the original backbone update at every layer. By induction on
$\ell$, all hidden states coincide with the frozen text-only model. Since the
output head is unchanged, the final logits and predictions also coincide.
\end{proof}

\begin{remark}
Proposition~\ref{prop:invariance} is an exact functional statement, not a
statistical one. It does not say that the model is robust to arbitrary
nonzero gates; it says that when the added modality path is truly zero, the
base model is preserved exactly.
\end{remark}

\section{Second-Order Interaction of Two Modalities}

To analyze composition locally, fix a hidden state $u \in \R^d$ at some fusion
site and let $\delta_a, \delta_v \in \R^d$ be the net residual perturbations
induced by audio and vision at that site. Let $q:\R^d \to \R$ be any smooth
scalar surrogate quantity of interest, for example a negative log-likelihood,
logit margin, or task score proxy.

\begin{proposition}[Second-order decomposition]
\label{prop:taylor}
Assume $q$ is three times continuously differentiable near $u$. Then
\begin{align}
q(u + \delta_a + \delta_v)
&=
q(u)
+ \inner{\nabla q(u)}{\delta_a + \delta_v}
+ \frac{1}{2}(\delta_a + \delta_v)^\top H_q(u)(\delta_a + \delta_v)
+ \bigO\!\left((\norm{\delta_a}+\norm{\delta_v})^3\right)
\nonumber\\
&=
\underbrace{q(u) + \inner{\nabla q(u)}{\delta_a}
+ \frac{1}{2}\delta_a^\top H_q(u)\delta_a}_{\text{audio-only terms}}
\underbrace{\inner{\nabla q(u)}{\delta_v}
+ \frac{1}{2}\delta_v^\top H_q(u)\delta_v}_{\text{vision-only terms}}
\nonumber\\
&\quad
+ \underbrace{\delta_a^\top H_q(u)\delta_v}_{\text{cross-term}}
+ \bigO\!\left((\norm{\delta_a}+\norm{\delta_v})^3\right).
\label{eq:taylor}
\end{align}
\end{proposition}

\begin{proof}
Apply the multivariate second-order Taylor expansion of $q$ at $u$ with
increment $\delta_a + \delta_v$:
\[
q(u+\eta)
=
q(u)
+ \inner{\nabla q(u)}{\eta}
+ \frac{1}{2}\eta^\top H_q(u)\eta
+ \bigO(\norm{\eta}^3),
\qquad
\eta := \delta_a + \delta_v.
\]
Expanding the quadratic term gives
\[
\frac{1}{2}(\delta_a+\delta_v)^\top H_q(u)(\delta_a+\delta_v)
=
\frac{1}{2}\delta_a^\top H_q(u)\delta_a
+ \frac{1}{2}\delta_v^\top H_q(u)\delta_v
+ \delta_a^\top H_q(u)\delta_v,
\]
since $H_q(u)$ is symmetric. Substituting this identity yields
Eq.~\eqref{eq:taylor}.
\end{proof}

\begin{remark}[Interpretation]
The cross-term $\delta_a^\top H_q(u)\delta_v$ is the leading source of local
non-additivity. If it is large and negative for a score-like $q$, joint
composition is worse than additive expectation; if it is positive, the two
modalities reinforce each other.
\end{remark}

\subsection{Gate scaling}

Write each perturbation as $\delta_m = g_m \hat \delta_m$, where $\hat \delta_m$
is the direction produced by the modality path and $g_m \ge 0$ is the gate.

\begin{corollary}[Quadratic scaling of cross-modal interference]
\label{cor:gate}
Under the notation above, the leading interaction term scales as
\[
\delta_a^\top H_q(u)\delta_v
=
g_a g_v \, \hat \delta_a^\top H_q(u)\hat \delta_v.
\]
Thus the cross-modal interaction is second-order in gate magnitude, while each
unimodal first-order term scales only linearly in its own gate.
\end{corollary}

\begin{proof}
Substitute $\delta_a = g_a \hat\delta_a$ and $\delta_v = g_v \hat\delta_v$
into the bilinear cross-term:
\[
\delta_a^\top H_q(u)\delta_v
=
(g_a \hat\delta_a)^\top H_q(u) (g_v \hat\delta_v)
=
g_a g_v \, \hat\delta_a^\top H_q(u)\hat\delta_v.
\]
\end{proof}

\begin{remark}
Corollary~\ref{cor:gate} explains why small initial gates can stabilize early
training: they suppress the interaction term faster than they suppress the
individual linear signals.
\end{remark}

\section{Synergy as an Empirical Proxy}

In experiments we do not usually optimize or report $q$ directly. Instead we
observe four task metrics:
\[
y_0 \quad \text{(text-only)}, \qquad
y_a \quad \text{(audio-only)}, \qquad
y_v \quad \text{(vision-only)}, \qquad
y_{av} \quad \text{(both)}.
\]

\begin{definition}[Synergy]
The synergy of the two-modality system is
\begin{equation}
\mathrm{syn}(y)
:=
y_{av} - (y_a + y_v - y_0).
\label{eq:synergy}
\end{equation}
\end{definition}

\begin{definition}[Gain vs best single]
The operational composition gain is
\begin{equation}
\Delta^{\mathrm{best}}
:=
y_{av} - \max(y_a, y_v).
\label{eq:gainbest}
\end{equation}
\end{definition}

\paragraph{Meaning.}
Equation~\eqref{eq:synergy} compares the observed composed score to the score
predicted by exact additivity. Equation~\eqref{eq:gainbest} answers the
practical question of whether composition beats the strongest unimodal system.

\begin{proposition}[Surrogate-level synergy identity]
\label{prop:synergy}
Let $q_0, q_a, q_v, q_{av}$ be the values of a smooth scalar surrogate under
text-only, audio-only, vision-only, and composed conditions, all linearized at
the same base hidden state $u$. Then
\begin{equation}
q_{av} - (q_a + q_v - q_0)
=
\delta_a^\top H_q(u)\delta_v
+ \bigO\!\left((\norm{\delta_a}+\norm{\delta_v})^3\right).
\label{eq:surrogate-synergy}
\end{equation}
\end{proposition}

\begin{proof}
Apply Proposition~\ref{prop:taylor} to $q_{av} = q(u+\delta_a+\delta_v)$ and
the corresponding second-order expansions for $q_a = q(u+\delta_a)$ and
$q_v = q(u+\delta_v)$. Subtracting $q_a + q_v - q_0$ cancels all constant,
gradient, and unimodal quadratic terms, leaving only the mixed quadratic term
plus third-order remainder.
\end{proof}

\begin{remark}
For discrete metrics such as exact-match accuracy, Eq.~\eqref{eq:surrogate-synergy}
is not an exact theorem. The experimental synergy in Eq.~\eqref{eq:synergy}
should be read as an empirical proxy for the same interference phenomenon, not
as the Hessian cross-term itself.
\end{remark}

\section{Why Same-Layer Fusion is More Fragile}

\subsection{Same-layer coupling}

Suppose both modalities inject at the same layer $\ell$:
\[
\tilde h_\ell = h_\ell + \delta_{a,\ell} + \delta_{v,\ell}.
\]
Then the next frozen block $f_\ell$ receives the \emph{sum} of both
perturbations simultaneously. By Proposition~\ref{prop:taylor}, the local
interaction at this layer is governed by the direct mixed term
\[
\delta_{a,\ell}^\top H_{f_\ell}(h_\ell)\delta_{v,\ell},
\]
where $H_{f_\ell}(h_\ell)$ denotes the appropriate second derivative of the
frozen layer map or of a scalar surrogate composed with that map.

\subsection{Disjoint-layer fusion}

Now suppose audio injects at layer $\ell_a$ and vision injects later at
$\ell_v > \ell_a$, with no shared injection at intermediate layers. Then
\[
\tilde h_{\ell_a} = h_{\ell_a} + \delta_{a,\ell_a},
\qquad
h_{\ell_v}
=
F_{\ell_a \to \ell_v}\!\left(h_{\ell_a} + \delta_{a,\ell_a}\right),
\]
where
\[
F_{\ell_a \to \ell_v}
:=
f_{\ell_v-1}\circ \cdots \circ f_{\ell_a}.
\]
Linearizing at $h_{\ell_a}$ gives
\[
h_{\ell_v}
\approx
F_{\ell_a \to \ell_v}(h_{\ell_a})
+ J_{\ell_a \to \ell_v}\,\delta_{a,\ell_a},
\]
where $J_{\ell_a \to \ell_v}$ is the Jacobian of the frozen backbone between
the two injection sites. Vision then injects
\[
\tilde h_{\ell_v}
=
h_{\ell_v} + \delta_{v,\ell_v}.
\]

\begin{proposition}[Direct cross-term elimination under disjoint layers]
\label{prop:disjoint}
If two modalities use disjoint layer sets, then there is no direct same-layer
mixed quadratic term of the form
$\delta_{a,\ell}^\top H_{f_\ell}(h_\ell)\delta_{v,\ell}$ at any fusion layer.
The leading interaction is instead Jacobian-mediated:
\begin{equation}
\delta_{v,\ell_v}^\top
H_{f_{\ell_v}}(h_{\ell_v})
\left(J_{\ell_a \to \ell_v}\,\delta_{a,\ell_a}\right),
\label{eq:indirect}
\end{equation}
or the symmetric analogue when vision precedes audio.
\end{proposition}

\begin{proof}
At any audio-only layer $\ell \in S_a \setminus S_v$, the perturbation entering
the frozen block is $h_\ell + \delta_{a,\ell}$, so no term involving
$\delta_{v,\ell}$ exists. Similarly, at any vision-only layer
$\ell \in S_v \setminus S_a$, the perturbation entering the frozen block is
$h_\ell + \delta_{v,\ell}$, so no direct same-layer mixed term exists there
either. The first point at which audio can influence the vision interaction is
through the transformed hidden state that arrives at a later vision layer,
which is given to first order by $J_{\ell_a \to \ell_v}\delta_{a,\ell_a}$.
Substituting this transformed perturbation into the mixed term at the later
vision layer yields Eq.~\eqref{eq:indirect}.
\end{proof}

\begin{corollary}[Interaction count]
With $M$ modalities sharing the same layer, the number of direct mixed pairs at
that layer is $\binom{M}{2}$. With disjoint layer sets, the number of direct
same-layer mixed pairs is zero.
\end{corollary}

\begin{proof}
In the shared-layer case, every unordered pair of active modality
perturbations contributes one mixed quadratic term. The number of unordered
pairs among $M$ modalities is $\binom{M}{2}$. Under disjoint layer sets, no
layer contains more than one modality perturbation, so no such same-layer pair
exists.
\end{proof}

\begin{remark}
Disjoint placement does \emph{not} eliminate all interaction. It only removes
the most direct one. The remaining interaction is mediated by the frozen
backbone Jacobian, which is typically weaker and easier to manage with gates
and training.
\end{remark}

\section{Diagnostics Used in Practice}

For a layer-level probe at layer $\ell$, define:
\[
y_{0,\ell}, \quad y_{a,\ell}, \quad y_{v,\ell}, \quad y_{av,\ell}.
\]

\begin{definition}[Layer additivity error]
The layer additivity error is
\begin{equation}
\epsilon_\ell
:=
\frac{\left|y_{0,\ell} - y_{a,\ell} - y_{v,\ell} + y_{av,\ell}\right|}
{\left|y_{a,\ell}-y_{0,\ell}\right| + \left|y_{v,\ell}-y_{0,\ell}\right| + \varepsilon},
\label{eq:additivity}
\end{equation}
with $\varepsilon > 0$ for numerical stability.
\end{definition}

\paragraph{Interpretation.}
Small $\epsilon_\ell$ means the layer behaves close to additive composition;
large $\epsilon_\ell$ indicates a high degree of non-additive interaction.

\paragraph{Why use both $\epsilon_\ell$ and $\Delta^{\mathrm{best}}_\ell$.}
The normalized error in Eq.~\eqref{eq:additivity} measures stability of
composition. The gain in Eq.~\eqref{eq:gainbest} measures usefulness of
composition. A strong fusion layer should ideally have both low
$\epsilon_\ell$ and high $\Delta^{\mathrm{best}}_\ell$.

\section{What These Results Support in the Paper}

The results above support four concrete claims:
\begin{enumerate}
\item \textbf{Functional invariance.} Added modalities do not alter the frozen
text hub when their residual paths are inactive
(Proposition~\ref{prop:invariance}).
\item \textbf{Interference has a specific local source.} The leading error of
joint composition is the mixed second-order term
(Proposition~\ref{prop:taylor}).
\item \textbf{Synergy is the right empirical summary.} The experimental synergy
metric is the observable proxy for this mixed interaction
(Proposition~\ref{prop:synergy}).
\item \textbf{Layer placement matters structurally.} Shared layers create direct
coupling; disjoint layers remove it and replace it with a weaker transformed
interaction (Proposition~\ref{prop:disjoint}).
\end{enumerate}

\paragraph{What these results do not claim.}
They do not prove that joint performance must exceed the best unimodal result,
that staggered layers always outperform shared layers on every backbone, or
that a global discrete metric is exactly equal to a local Hessian quantity.
Those are empirical questions.

\end{document}
